% from ega2-indexes-and-contents-fr.tex
\oldpage[II]{213}
\markboth{TABLE DES MATIÈRES}{TABLE DES MATIÈRES}
\phantomsection
\addcontentsline{toc}{section}{Table des matières originale}
\section*{Table des matières}
\thispagestyle{plain}

\begingroup
\small
\renewcommand{\arraystretch}{1.04}
\begin{longtable}{@{}p{0.09\textwidth}p{0.69\textwidth}r@{}}
\textbf{Unité} & \textbf{Titre} & \textbf{Page originale}\\
\hline
\textbf{II} & \textbf{Étude globale élémentaire de quelques classes de morphismes} & \textbf{5}\\
\hyperref[section:II.1-fr]{\textsection1} & Morphismes affines & 5\\
\hyperref[subsection:II.1.1-fr]{1.1} & $S$-préschémas et $\mathcal{O}_S$-Algèbres & 5\\
\hyperref[subsection:II.1.2-fr]{1.2} & Préschémas affines sur un préschéma & 6\\
\hyperref[subsection:II.1.3-fr]{1.3} & Préschéma affine au-dessus de $S$ associé à une $\mathcal{O}_S$-Algèbre & 8\\
\hyperref[subsection:II.1.4-fr]{1.4} & Faisceaux quasi-cohérents sur un préschéma affine au-dessus de $S$ & 9\\
\hyperref[II.1.5-fr]{1.5} & Changement du préschéma de base & 12\\
\hyperref[II.1.6-fr]{1.6} & Morphismes affines & 14\\
\hyperref[II.1.7-fr]{1.7} & Fibré vectoriel associé à un faisceau de modules & 14\\
\hyperref[section:II.2-fr]{\textsection2} & Spectres premiers homogènes & 19\\
\hyperref[subsection:II.2.1-fr]{2.1} & Généralités sur les anneaux et modules gradués & 19\\
\hyperref[subsection:II.2.2-fr]{2.2} & Anneaux de fractions d'un anneau gradué & 23\\
\hyperref[subsection:II.2.3-fr]{2.3} & Spectre premier homogène d'un anneau gradué & 25\\
\hyperref[subsection:II.2.4-fr]{2.4} & La structure de schéma sur $\Proj(S)$ & 28\\
\hyperref[subsection:II.2.5-fr]{2.5} & Faisceau associé à un module gradué & 30\\
\hyperref[subsection:II.2.6-fr]{2.6} & $S$-module gradué associé à un faisceau sur $\Proj(S)$ & 36\\
\hyperref[subsection:II.2.7-fr]{2.7} & Conditions de finitude & 38\\
\hyperref[subsection:II.2.8-fr]{2.8} & Comportements fonctoriels & 41\\
\hyperref[subsection:II.2.9-fr]{2.9} & Sous-préschémas fermés d'un schéma $\Proj(S)$ & 48\\
\hyperref[section:II.3-fr]{\textsection3} & Spectre homogène d'un faisceau d'algèbres graduées & 49\\
\hyperref[subsection:II.3.1-fr]{3.1} & Spectre homogène d'une $\mathcal{O}_Y$-Algèbre graduée quasi-cohérente & 49\\
\hyperref[subsection:II.3.2-fr]{3.2} & Faisceau sur $\Proj(\mathcal{S})$ associé à un $\mathcal{S}$-Module gradué & 54\\
\hyperref[subsection:II.3.3-fr]{3.3} & $\mathcal{S}$-Module gradué associé à un faisceau sur $\Proj(\mathcal{S})$ & 56\\
\hyperref[subsection:II.3.4-fr]{3.4} & Conditions de finitude & 59\\
\hyperref[subsection:II.3.5-fr]{3.5} & Comportements fonctoriels & 61\\
\hyperref[subsection:II.3.6-fr]{3.6} & Sous-préschémas fermés d'un préschéma $\Proj(\mathcal{S})$ & 64\\
\hyperref[subsection:II.3.7-fr]{3.7} & Morphismes d'un préschéma dans un spectre homogène & 65\\
\hyperref[subsection:II.3.8-fr]{3.8} & Critères d'immersion dans un spectre homogène & 69\\
\end{longtable}
